\chapter{Lista modeli 3D}
W ramach pracy powstały następujące modele 3D:
\begin{itemize}
    \item \textbf{VibrometerPositioningTable.f3d} -- Główne złożenie całego stołu do pozycjonowania wibrometru.
    \item Dolna, stacjonarna część stołu:
    \begin{itemize}
        \item \textbf{StaticBase.f3d} -- Cześć podtrzymująca dolne łożysko, do której przymocowany jest silnik.
        \item  \textbf{DownTable.f3b} -- Część która pozwala na montaż na stojaku mikrofonowym z gwintem \textbf{M6}.
        \item  \textbf{BaseSpacer.f3b} -- Część łącząca \textbf{StaticBase} oraz \textbf{DownTable}, dająca przestrzeń na montaż silnika.
    \end{itemize}
    \item Środkowa część stołu, poruszająca się w jednej osi.
    \begin{itemize}
        \item \textbf{PartialMovementTable.f3d} -- Część stołu, która porusza się dookoła osi z. Dodatkowo zawiera otwory, aby przymocować statyczną część drugiego silnika.
        \item \textbf{HubSpacer2.f3b} \& \textbf{HubSpcer5.f3b} -- Modele, które razem z hubem montażowym Pololu 1083 tworzą sprzęgiełko łączące \textbf{PartialMovementTable} wraz z dolnym silnikiem.
    \end{itemize}
    \item Górna część stołu, która wraz z wibrometrem obraca się dookoła obu osi.
    \begin{itemize}
        \item \textbf{Holder1.f3d} \& \textbf{Holder2.f3d} -- Dwie części, jedna montowalna do wału drugiego silnika, druga bezpośrednio na łożysko kulkowe.
        \item \textbf{PlywoodPlate.f3d} -- Cześć łącząca holdery do właściwej głowicy wibrometru.
    \end{itemize}
    
\end{itemize}